\documentclass[10pt]{article}
\usepackage[margin=20mm]{geometry}
\usepackage{amsmath,amssymb,fontspec,unicode-math,xurl,hyperref}
\IfFontExistsTF{DejaVu Serif}{\setmainfont{DejaVu Serif}}{\setmainfont{Latin Modern Roman}}
\setmathfont{latinmodern-math.otf}
\setlength{\parindent}{0pt}\setlength{\parskip}{6pt}\sloppy
\begin{document}
Abstract: The published oracle definition makes the asserted stochastic lower bound false at hitting-time parameter one. The rare-state construction also fails to preserve that parameter. Independently, MaC-SAGE’s displayed recursion does not maintain the weighted gradient used in its convergence proof. These are major defects in core guarantees. A separate minor defect concerns the motivating temporal-difference example. The analysis does not establish that every possible corrected theorem is false. Adjudication is pending.\par
Author: AI-assisted technical review\par
Date:\par
Title: Technical Verification of Markov-Sampling Lower Bounds\par
\# Identity and inspected scope\par
Zhenyu Sun and Ermin Wei, Improved Lower Bounds for First-order Stochastic Non-convex Optimization under Markov Sampling, ICML 2025, PMLR 267, 57754–57772; identifier I25-43690 (0321), OpenReview `vBYRG3p21R`. The reviewed version is the frozen 19-page publisher PDF. Page numbers below are PDF pages, which equal its printed pages. All main-paper and appendix text was read. Claim-level checks concentrated on Definitions 2.1 and (6), Theorem 3.1, the construction in Appendix A.4, Algorithm 1, the decomposition in Theorem B.5, and Section 2.1. Figures 1–2 and the relevant equation pages were visually inspected. There are no reproduced experiments. No released code was available in the packet; the separate OpenReview attachment inventory was inaccessible. The remaining lemmas have been read but not independently proved in full.\par
\# Definitions that matter\par
Equation reference: eq:hit. The oracle returns \(g(x;s)=\nabla f(x;s)\) with stationary expectation \(\nabla F(x)\) and stationary variance at most \(\sigma^2\). The function class has \(L\)-Lipschitz gradient and \(F(0)-\inf F\leq\Delta\). The chain class in (6) imposes \[
 \tau_{\rm hit}=\max_{v,w}E_v\inf\{t\geq1:s_t=w\}\leq\tau.\] In particular, the maximum includes return times (\(v=w\)), not just mixing times. Zero-respecting algorithms can use gradient descent from zero, including selection of a previously queried iterate.\par
\# Substantive findings\par
Problem 1 (Theorem 3.1 is incompatible with its hitting-time class).  This is a major defect in the principal stochastic lower bound.\par
\#\# Source claim and decisive check\par
Equation reference: eq:claimed. Theorem 3.1 (p. 5) asserts, for positive parameters including \(\tau=1\), \[
 N_s^\epsilon=\Omega\left(\frac{\tau\sigma^2}{\epsilon^2}
 +\frac{\tau\sigma^2}{\epsilon^4}
 \min\{c_1\sigma^2,c_2L\Delta\}\right),\qquad c_1,c_2>0.\] Remark 3.2 explicitly discusses \(\tau=1\). Under [eq:hit] (\#eq:hit), however, this class has just one-state chains. Indeed, any nonnegative integer hitting time is at least one. If its expectation is at most one, it is identically one. Thus for fixed \(v\), \(P(v,w)=1\) for every \(w\), impossible for two distinct states. A singleton chain has stationary probability one, so stationary unbiasedness implies \(g(x;s)=\nabla F(x)\) exactly. Its variance is zero, which satisfies the stated upper bound for every \(\sigma>0\).\par
For every admissible \(F\), deterministic gradient descent with step \(1/L\) satisfies the smoothness descent inequality \[F(x_{t+1})\leq F(x_t)-\frac{1}{2L}\|\nabla F(x_t)\|^2,
 \qquad
 \min_{t<T}\|\nabla F(x_t)\|^2\leq\frac{2L\Delta}{T}.\] It is zero-respecting and uses one query per step. Selecting the queried iterate with smallest observed gradient therefore gives \(N_s^\epsilon=O(L\Delta\epsilon^{-2})\) at \(\tau=1\). Fix \(L=\Delta=\sigma=1\) and let \(\epsilon\downarrow0\). This contradicts the positive \(\epsilon^{-4}\) term in [eq:claimed] (\#eq:claimed). No special choice of objective, variance lower bound, or hidden regularity assumption is needed: the upper bound holds for the entire displayed class at this parameter.\par
\#\# Why the construction does not resolve this\par
Appendix A.4 (pp. 13–15, Fig. 2) splits \(v^*\) into states with stationary masses \(q\pi_{v^*}\) and \((1-q)\pi_{v^*}\), and chooses \(q^{-1}=1+\sigma^2/(8a_2\epsilon^2)\). For a recurrent irreducible chain, the expected positive return time to a state of mass \(p\) is \(1/p\): successive return cycles have mean length \(m\), and their stationary visit frequency is \(1/m=p\). Hence the new chain obeys \[\widetilde\tau_{\rm hit}\geq\frac{1}{q\pi_{v^*}}\geq\frac1q.\] For fixed \(\tau\), sufficiently small \(\epsilon\) puts it outside (6). For example, \(q=0.01\) forces a return time at least \(100\), inconsistent with \(\tau=10\). This necessary condition does not depend on how the unpictured transition probabilities are completed. More generally, a finite uniform return-time bound forces every stationary mass to be at least \(1/\tau\), so the irreducible state space has at most \(\lfloor\tau\rfloor\) states. Allowing a countable state space does not remove the obstruction.\par
\#\# Counterevidence, impact, and repair scope\par
An i.i.d. chain on several states has mixing time one, but hitting time \(\max_w1/\pi_w\), not one. Interpreting \(\tau\) as mixing time could avoid the singleton argument, but contradicts Definition 2.1 and (6), and requires a new proof and parameter accounting. The stated core theorem is false in its explicit parameter range; the supplied construction also fails for small \(q\) at fixed hitting bound. This does not refute stochastic lower bounds under other oracle classes. No repaired lower bound is certified here. The separate weighted-gradient lower bound in Lemma A.1 also requires normalization scrutiny; it is not needed for this verdict.\par
Problem 2 (Algorithm 1 does not maintain the gradient sum used in Theorem B.5).  This is a major invalid derivation in the core MaC-SAGE guarantee.\par
\#\# Claim and fully admissible calculation\par
Equation reference: eq:actual. On p. 6, line 5 sets \[
 G_t=G_{t-1}-y_{t-1}^i h_{t-1}^i+y_t^i\nabla f_i(x_t),
 \quad s_t=i.\] The prose then claims \(G_t=\sum_jy_t^jh_t^j\). This equality is used to bound \(\|\nabla F(x_t)-G_t\|^2\) in Theorem B.5 (p. 18). Consider two states sampled independently with probabilities \(1/2,1/2\), \(f_1(x)=x^2/2+x\) and \(f_2(x)=x^2/2-x\). Then \(F(x)=x^2/2\), \(L=1\), the initial gap at \(x_0=0\) is zero (thus at most \(\Delta=1\)), and \(|\nabla f_i(x)-\nabla F(x)|=1\) for every \(x\). Assumption 4.1 holds with \(\sigma=1\) and the hitting bound is \(\tau=2\).\par
The path \(s_0=1,s_1=2\) has probability \(1/4\). With the prescribed zero initial tables and any positive \(\gamma_0\) (take \(1/8\)), \[\begin{aligned}
 y_0&=(1,0),&h_0&=(1,0),&G_0&=1,&x_1&=-1/8,\\
 y_1&=(1/2,1/2),&h_1&=(1,-9/8),&G_1&=7/16.\end{aligned}\] But \(\sum_jy_1^jh_1^j=1/2-9/16=-1/16\). The discrepancy is \(1/2\); it occurs because the unvisited state’s weight decays but its contribution in [eq:actual] (\#eq:actual) does not. The discrepancy is independent of the selected step size. Exact rational arithmetic in the accompanying independent script reproduces both values.\par
\#\# Repair and bounded consequence\par
Equation reference: eq:repair. For the intended invariant \(G_t=\sum_jy_t^jh_t^j\), a correct update is \[
 G_t=\frac{t}{t+1}G_{t-1}
 -\frac{t}{t+1}y_{t-1}^i h_{t-1}^i
 +y_t^i\nabla f_i(x_t).\] To verify it, use \(y_t^j=t y_{t-1}^j/(t+1)\) and \(h_t^j=h_{t-1}^j\) for \(j\ne i\), then replace just the \(i\)th term. The zero initialization also makes this valid at \(t=0\). This repairs the invariant only. It does not validate the rest of Appendix B, including the application of a fixed-function moment bound to iterate-dependent gradients. Alternatively, one could analyze the printed algorithm with its extra accumulated error, but that analysis is absent. The finite-path example refutes the claimed equivalence and its use in the proof; it does not by itself prove asymptotic divergence or rule out the stated rate after repair. Because the defect invalidates the derivation of the proposed algorithm’s central guarantee, its impact is major.\par
Problem 3 (The TD line integral is not a potential for the TD vector field).  This is a minor defect in the motivating application, without a shown dependency in the abstract lower-bound theorems.\par
Section 2.1 (p. 3) defines \(f(\theta;\bar s)=\int_0^1g(\theta_0+u(\theta-\theta_0);s,s')^T
(\theta-\theta_0)\,du\) and concludes \(\nabla f=g\). Take \(\theta_0=0\), \(\gamma=1/2\), reward zero, \(\phi(s)=e_1\), \(\phi(s')=e_2\). These features can be assigned to the two states of the preceding i.i.d. chain. The true discounted value is identically zero, so the assumed exact linear representation holds with \(\theta^*=0\). The transition \(s\to s'\) has positive probability. For \(\theta=(x,y)\) the specified formulas give \[g(\theta;s,s')=(x-y/2,0)^T,\quad
 f=x^2/2-xy/4,\quad
 \nabla f=(x-y/4,-x/4)^T.\] At \((0,1)\) these are \((-1/2,0)^T\) and \((-1/4,0)^T\). The Jacobian of \(g\) is nonsymmetric, so no differentiable scalar potential with this gradient exists on \(\mathbb R^2\). A conservative vector field would make the line-integral argument valid, but that condition is neither assumed nor supplied by linear value representation. Even if the expected TD field is conservative in a special chain, that does not establish the asserted per-sample identity. The proposed reformulation of ordinary TD updates is therefore invalid in general.\par
\# Overall impact and state\par
Theorem 3.1 fails as stated, and the MaC-SAGE convergence proof does not apply to its displayed recurrence. These concern two central contributions, irrespective of whether related results can be recovered. Theorem 4.2 and all remaining proof steps are not certified by this review. No experimental scores are challenged, since the paper is theoretical. No later correction was present in the frozen packet; a later-version/erratum comparison was not performed and author adoption is unknown.\par
Dimension | State\par
Supported overall assessment | Major\par
P1 / P2 / P3 | Major / Major / Minor\par
Repair verification | P2 invariant only, equation [eq:repair] (\#eq:repair)\par
Scientific adjudication | Pending\par
Human review / author adoption | None recorded / unknown\par
Remaining work | New lower-bound and convergence proofs\par
1 Z. Sun and E. Wei, Improved Lower Bounds for First-order Stochastic Non-convex Optimization under Markov Sampling, PMLR 267 (2025), 57754–57772. Frozen publisher version, 19 pp. https://proceedings.mlr.press/v267/sun25u.html.\par
\# Evidence and reproducibility\par
The source SHA-256 is `b408c2e4f735c9560f2a48e38e0b171f60f5f3060a1398886756e04516866552`. The matching verification directory contains `source\_utf8.txt`, `source\_check.json`, the inspected source-page renders, `independent\_checks.py` and its actual JSON output, and the coverage, findings, build, visual-inspection and provenance records. The script uses Python 3.13.9 and SymPy 1.14.0 for exact TD and recurrence arithmetic; it also reports the diagnostic \(\|\nabla(\tfrac12h_1+\tfrac12h_2)(0)\|=\sqrt e/2\) for the weighted Lemma A.1 formula. That diagnostic is not relied on as a full reconstruction of Appendix A.1’s unspecified transition matrix. No author program was executed. The source reading, proofs, calculations, and visual checks are by the initial AI reviewer selected as GPT-6 Astra with Max effort; they are not independent acceptance or human review. Full hash-bound execution details are external to the scientific narrative.\par
\end{document}
